Cardiovascular diseases (CVDs) remain the leading cause of mortality globally, claiming approximately 17.9 million lives annually. Traditional diagnostic tools, such as 12-lead electrocardiography (ECG) and wearable Holter monitors, require specialized instrumentation, conductive gels, and adhesive electrodes, while consumer smartwatches still require \$300+ dedicated wearable hardware and rely on optical wrist photoplethysmography (PPG) without capturing direct chest valve mechanics. To address these barriers, this senior capstone project presents an end-to-end, zero-external-hardware mobile cardiovascular screening, binary arrhythmia classification, and controlled arrhythmia simulation platform executed by a two-member engineering team. The system captures mechanical sternal seismocardiography (SCG) micro-vibrations along the dorso-ventral ($Z$-axis) direction using the built-in 3-axis micro-electromechanical system (MEMS) accelerometer of a commercial smartphone placed flat on the subject's lower sternum in a supine resting posture.

To eliminate non-biological noise and motion artifacts, an automated digital signal processing pipeline was engineered, incorporating a zero-phase 4th-order Butterworth bandpass filter (0.5--40.0 Hz) and an absolute physiological energy gate ($0.005\text{ m/s}^2 \le \sigma_{raw} \le 0.50\text{ m/s}^2$) paired with real-time stillness guidance. Peak-to-peak interval (PPI) tachograms are derived via a lightweight 200 ms moving-RMS energy envelope and adaptive sliding thresholding. Extracted pulse rate variability (PRV) features—including time-domain (Mean HR, SDNN, RMSSD, pNN50) and frequency-domain (Welch spectral LF, HF, and LF/HF) metrics—feed into explainable supervised machine learning classifiers focused strictly on a single abnormality: binary Cardiac Arrhythmia classification (Normal sinus rhythm [$N$] versus Abnormal / Arrhythmia [$S, V, F, Q$]).

Evaluated on the PhysioNet MIT-BIH Arrhythmia Database, XGBoost achieved an overall binary classification accuracy of 89.4\%, a sensitivity (recall) of 88.9\%, a specificity of 89.8\%, and an F1-score of 0.887. Systolic peak detection yielded a pulse rate Mean Absolute Error of 2.14 bpm against gold-standard ECG, with Bland--Altman 95\% limits of agreement spanning $[-4.18, +4.82]\text{ bpm}$. The total edge execution latency was 56.1 ms, comfortably satisfying soft real-time constraints ($<500\text{ ms}$). Finally, a coupled nonlinear Van der Pol oscillator (RK45) serves as a controlled arrhythmia simulation tool to generate synthetic arrhythmic SCG test signals, bridging the scarcity of real pathological smartphone SCG recordings during end-to-end verification. This work confirms the viability of smartphone sternal SCG for accessible, zero-hardware cardiac arrhythmia screening.
